\section{Random Forest (Thomas Wobak)}
This part was worked on by Thomas Wobak (Matnr: 12203494).
The goal of this implementation was to build a Random Forest (RF) classifier for the project. With the preprocessing already being done, it was just a work of importing it and starting the work.
\nextparagraph
The core concept of the RF approach was to not only build a standars RF classifier, but also to improve it by correcting the individual trees and dynamically correcting predictions. In general, the RF implementation was split into four phases:
\nextparagraph
\begin{enumerate}
    \item \textbf{Establish a Baseline} the objective of the first step was to set up a foundation to measure the impact of the upcoming model enhancements against. For this step I loaded the preprocessed data and trained a standard RandomForestClassifier for both the binary task and the multiclass task.
    \item \textbf{Implementing the Enhanced Random Forest} After obtaining the baseline stats, the next step was to optimize the RF ensemble by ensuring maximum tree diversity and individual accuracy. This was done by dissecting the forest to evaluate the individual decision trees using the Area Under Curve (AUC) metric. This was done by using the Jaccard index of the split features between trees. To prevent overfitting or collapsing, the tree depth was capped at 20 and a similarity threshold was applied.
    \item \textbf{The Result Correction Mechanism} After obtaining the enhanced RF, the third step applied a secondary validation layer, specifically to lower the False Alarm Rate (FAR) and correct predictions that underperformed. This was done via the ResultCorrector class. For the binary model a probability threshold was introduced which required a 65\% confidence level for the model, while for the multiclass a confidence-based tie-breaker was applied. If the RF confidence fell below 55\% a specific feature rule (e.g. sload levels for DoS, dur and spkts for Reconnaissance) was applied to override the prediction. This prevented aggressive label bias (during testing it would predict almost 80\% of all traffic as Reconnasaince before this) while correcting uncertain labels.
    \item \textbf{Comprehensive Evaluation and Visualization} The final step was to validate the enhanced model against the projects performance goals and visualize the results. This showed that the binary classifier performed great, achieving an accuracy score of 0.9188 with an FAR of 0.1107. For the multiclass model the performance was worse, but still performed well, achieving an accuracy of 0.6859.
\end{enumerate}
The final results can be seen in the following table:
\begin{table}[htbp]
\centering
\caption{Performance Evaluation of the Baseline and Corrected Random Forest Models}
\label{tab:rf_metrics}
\begin{tabular}{l@{\hspace{0.3cm}}lcccccc}
\hline
\textbf{Task} & \textbf{Model Version} & \textbf{Accuracy} & \textbf{Precision} & \textbf{Recall (DR)} & \textbf{F1-score} & \textbf{MCC} & \textbf{FPR (FAR)} \\ \hline
Binary & Baseline & 0.8979 & 0.8591 & 0.9743 & 0.9131 & 0.7997 & 0.1958 \\
Binary & Corrected & 0.9188 & 0.9126 & 0.9429 & 0.9275 & 0.8359 & 0.1107 \\ \hline
Multiclass & Baseline & 0.6862 & 0.8511 & 0.6862 & 0.7379 & 0.6233 & - \\
Multiclass & Corrected & 0.6859 & 0.8505 & 0.6859 & 0.7374 & 0.6229 & - \\ \hline
\end{tabular}
\end{table}
Unfortunately, the multiclass model did not improve much during the correction phase, but the binary model showed great metric improvements. The overall performance for the binary classifier still beats the goals of the project, with the goal performance of accuracy $> 0.9$ and FAR $< 0.15$ being achieved.
I am not sure why the multiclass model did not improve, but I suspect that it is because of the conservative correction logic in the ResultCorrector. During the first implementation the multiclass predictor would label over 80\% of all traffic as Reconnasaince attacks, so during the correction process I updated the logic to only override a prediction in the Random Forest if the confidence was very low and the strict feature thresholds are met. Because it no longer "blindly" overrides predictions, it alters very few records overall, leaving the baseline accuracy largely unchanged.
I tried altering the ResultCorrector logic to be more aggressive, but this led to the same result as previously. It might also be a limitation of the algorithm itself.
Since Random Forests are ensemble models that rely on majority voting, the result of the pruning during phase two might have created a RF where all trees are structurally sound but share the same fundamental limitation. So while the RF was an exceptional tool for the binary threat detection, it proved insufficient for precise attack categorization. I believe this could still be achieved but would require much deeper, attack-specific feature engineering.
\nextparagraph
There are also some design decisions that should be discussed for the RF classification. The first decision I would like to discuss further is the capping of the forest depth in favor of more diversity. The \texttt{max\_depth} of the decision trees was intentionally capped at 20, and the \texttt{max\_samples} was restricted to 0.8 of the training data. During the first phase allowing the trees to grow infinitely deep resulted in the multiclass model collapsing into a single tree during the pruning phase. I believe this was because the SMOTE algorithm perfectly balanced all nine classes, fully grown trees were memorizing the dataset and using the same features. By capping the depth the trees were forced to generalize, ensuring the pruned ensemble remained diverse.
Another design decision was pruning via AUC and Jaccard similarity. Instead of using a standard RF, I implemented a custom pruning function that evaluated individual trees based on their AUC and discarded redundant trees based on Jaccard similarity of their split features. This directly aligns with the methodology proposed by \cite{wu2022intrusion}.
By stripping out trees that were highly similar to each other but had poorer classification performance, the structural redundancy of the forest was reduced. This was done to optimize the final majority-voting process.
\nextparagraph
Overall, the RF implementation was a success, especially for the binary classification task. The multiclass model did not perform as well, but it still achieved a respectable accuracy of 0.6859. The RF implementation was a great learning experience, and I believe that with more time and deeper feature engineering, the multiclass model could be improved to the same level of the binary classification. 